\chapter{Diskussion}
\label{ch:diskussion}

Hier schreibe ich was, dass ich noch überarbeiten muss. Und

srzhearearzh

Die berechneten Flächen der kryogenen Wärmeübertrager liegen in einer
Größenordnung von etwa \(2{,}7 \cdot 10^3\) bis
\(1{,}36 \cdot 10^4\,\mathrm{m^2}\). Diese Größenordnung ist für kryogene
Luftzerlegungsanlagen plausibel und steht im Einklang mit den von Tesch für
vergleichbare Anlagenkonzepte angegebenen flächenbasierten Kostenparametern, bei
denen Wärmeübertragerflächen ebenfalls im Bereich mehrerer tausend bis über
\(10^4\,\mathrm{m^2}\) auftreten. Die Flächen der Zwischenkühler fallen aufgrund
der größeren Temperaturdifferenzen und der warmen Prozessbedingungen deutlich
kleiner aus. Die berechneten Werte sind als äquivalente Wärmeübertragerflächen zu
verstehen und dienen primär der flächenbasierten Kostenabschätzung; sie ersetzen
keine detaillierte Apparateauslegung.